\section{Related Work}
% ──────────────────────────────────────────────────────────────────────────────

Efforts to taxonomise XAI methods have produced a substantial but fragmented
body of work. Arrieta et al.~\cite{Arrieta2020XAI} established the foundational
distinction between interpretable-by-design models and post-hoc explanations,
alongside key challenges for responsible AI. Schwalbe and
Finzel~\cite{schwalbe2024comprehensive} synthesised over fifty surveys into a
framework encompassing problem definitions, explanator properties, and evaluation
metrics, demonstrating that properties such as faithfulness and stability are
widely referenced but inconsistently operationalised. Speith~\cite{speith2022review}
identified a structural source of this inconsistency: taxonomies grounded in
mechanisms versus those grounded in outputs classify the same methods
differently, limiting cross-study comparability. Collectively, these works
underscore the absence of a shared vocabulary for organising evaluation
criteria, particularly in high-stakes domains such as clinical AI.

Medical XAI taxonomies begin to address this gap but remain incomplete.
Jin et al.~\cite{jin2023guidelines} distinguish technical faithfulness from
clinical usefulness, providing a domain-specific framework that separates
algorithmic fidelity from practical utility.
Holzinger et al.~\cite{holzinger2022information} emphasise robustness under
distributional shift as an unresolved concern for clinical deployment, while Tjoa and
Guan~\cite{Tjoa2021MedicalXAI} survey XAI methods broadly without formalising
evaluation criteria. Varghese et al.~\cite{Varghese2025OnDE} rank XAI metrics
using Borda count and multiple correlation techniques,
identifying faithfulness and sensitivity as primary evaluative axes, though
without cross-domain validation. Reviews in radiology and nuclear
medicine~\cite{groen2022systematic,de2023explainable,saw2025current} further
demonstrate that evaluation metrics are inconsistently reported and rarely
subjected to comparative analysis. Reyes et al.~\cite{reyes2020interpretability}
argue that explanations must align with clinical reasoning, a requirement
frequently absent from benchmark-driven evaluations.

The metrics literature exposes the practical consequences of this
fragmentation. Rosenfeld~\cite{rosenfeld2021better} demonstrates that most
evaluations prioritise algorithmic fidelity while neglecting human-centred
dimensions. Kadir et al.~\cite{kadir2023evaluation} confirm that sensitivity,
faithfulness, and robustness dominate existing evaluation practice, with limited
attention to human factors. Phuong et al.~\cite{phuong2023benchmarking} and
Klein et al.~\cite{Klein2024NavigatingTM} provide evidence that, without
structured taxonomies, metric selection remains ad hoc and results
non-comparable across studies. Proposed evaluation frameworks such as that of
Mirzaei et al.~\cite{mirzaei2023xai} offer selection criteria but lack
grounding in domain-specific medical contexts. Complementary
analyses~\cite{coroama2022evaluation,boggust2023saliency} reinforce that the
absence of a shared evaluation vocabulary continues to impede meaningful
synthesis.

Human-centred evaluation has received increasing attention in recent literature.
Chen et al.~\cite{chen2022explainable} identify the systematic absence of
human-centred evaluation components in medical XAI pipelines, while Bauer and
Michalowski~\cite{bauer2025human} highlight cognitive load, trust calibration,
and decision accuracy as key constructs requiring operationalisation. Dey et
al.~\cite{dey2022human} document continued methodological fragmentation in
human-centred XAI evaluation, and M{\"u}ller et
al.~\cite{muller2022explainability} emphasise that regulatory demands for both
explainability and causability extend beyond standard faithfulness metrics,
necessitating evaluation approaches that span technical and user-facing
dimensions.

Recent surveys~\cite{houssein2025explainable,zhang2026comprehensive,Ahmed2026XAI}
confirm that, despite extensive method-level research, standardised evaluation
frameworks remain underdeveloped. In particular, Ahmed et
al.~\cite{Ahmed2026XAI} identify metric standardisation as a major open
challenge across medical imaging studies.

What remains absent from the literature is a unified taxonomy for organising
XAI evaluation across medical sub-domains, grounded in empirical usage
patterns rather than purely conceptual design. No existing framework jointly
spans oncology, cardiology, neuroimaging, radiology, and clinical decision
support while integrating mathematical faithfulness, system robustness, and
human-centred trust into a single evaluative structure. The present study
addresses this gap by conducting a systematic co-occurrence analysis and
constructing an empirically validated three-layer taxonomy.
